\section{Error Handling}
\label{sec:error-handling}

In case you encounter errors while using TOX, the following error codes may
help you identify the issue. The error codes are grouped into categories based
on their nature, such as I/O errors, format validation errors, memory errors, and
internal errors.

\begin{longtable}
    {>{\ttfamily}p{3cm} p{10cm}}
    \caption{TOX Error Codes Reference}
    \\ \hline \textbf{Error Code} & \textbf{Description} \\ \hline \endfirsthead

    \multicolumn{2}{c}%
    {{\bfseries \tablename\ \thetable{} -- continued from previous page}} \\ \hline
    \textbf{Error Code} & \textbf{Description} \\ \hline \endhead

    \hline \multicolumn{2}{r}{{Continued on next page}} \\ \endfoot

    \hline \endlastfoot

    0 & Success \\[6pt]

    \multicolumn{2}{l}{\textbf{1xx: I/O / File Reading Errors}} \\[3pt] 
    101 & Could not open file. \\ 
    102 & Could not read magic number. \\ 
    103 & Could not read array type code. \\ 
    104 & Could not read number of dimensions. \\ 
    105 & Could not read array dimensions. \\ 
    106 & Could not read character length. \\ 
    107 & Could not read array data. \\ 
    112 & Could not write magic number \\ 
    113 & Could not write array type code\\ 
    114 & Could not write number of dimensions \\
    115 & Could not write dimensions \\ 
    116 & Could not write character length \\
    117 & Could not write array data \\[6pt]

    \multicolumn{2}{l}{\textbf{2xx: Format / Input Validation Errors}} \\[3pt]
    200 & Invalid format detected (e.g., magic mismatch or corrupt format). \\ 
    201 & Invalid input provided. \\ 
    202 & Empty input arrays provided (e.g., zero-length arrays). \\ 
    203 & Dimension mismatch detected (shape inconsistencies). \\
    204 & NaN or Inf found in input data where not allowed. \\ 
    205 & Unsupported data type encountered (e.g., complex types). \\ 
    206 & Array size mismatch \\
    208 & Array index out of bounds \\ 
    209 & Division by zero detected \\[6pt]

    \multicolumn{2}{l}{\textbf{3xx: Memory Errors}} \\[3pt] 
    301 & Memory allocation failed. \\ 
    302 & Null pointer reference encountered. \\[6pt]

    \multicolumn{2}{l}{\textbf{5xxx: Fortran Runtime / Unit State Errors}} \\[3pt] 
    5002 & Fortran runtime error: unit not open or not connected. \\[6pt]

    \multicolumn{2}{l}{\textbf{9xxx: Internal / Unknown Errors}} \\[3pt] 9001 & Internal
    error: unexpected state or invariant violated. \\ 9999 & Unknown error (fallback
    for unexpected errors). \\
\end{longtable}